\documentclass[10pt,letterpaper]{article}
\usepackage[T1]{fontenc}
\usepackage[utf8]{inputenc}
\usepackage[sfdefault]{sourcesanspro}
\usepackage{inconsolata}
\usepackage[margin=0.84in,headheight=15pt]{geometry}
\usepackage{microtype}
\usepackage{xcolor}
\usepackage{booktabs,longtable,tabularx,array}
\usepackage{enumitem}
\usepackage{fancyhdr,lastpage}
\usepackage[most]{tcolorbox}
\usepackage{listings}
\usepackage{titlesec}
\usepackage{hyperref}
\usepackage{amsmath,amssymb}
\usepackage{multicol}
\definecolor{navy}{HTML}{15324C}
\definecolor{teal}{HTML}{087E83}
\definecolor{slate}{HTML}{536879}
\definecolor{paper}{HTML}{F2F6F8}
\definecolor{warm}{HTML}{E8F4F3}
\hypersetup{colorlinks=true,urlcolor=teal,linkcolor=navy,pdfauthor={Software Engineering Reference Suite},pdfsubject={MISRA C++:2023 engineering guidance}}
\titleformat{\section}{\Large\bfseries\color{navy}}{\thesection}{0.7em}{}
\titleformat{\subsection}{\large\bfseries\color{teal}}{\thesubsection}{0.6em}{}
\titleformat{\subsubsection}{\normalsize\bfseries\color{navy}}{\thesubsubsection}{0.5em}{}
\setlist[itemize]{leftmargin=*,topsep=2pt,itemsep=2pt}
\setlist[enumerate]{leftmargin=*,topsep=2pt,itemsep=3pt}
\setlength{\parskip}{5pt}
\setlength{\parindent}{0pt}
\setlength{\emergencystretch}{2em}
\pagestyle{fancy}
\fancyhf{}
\fancyhead[L]{\small\color{slate}MISRA C++:2023 Engineering Suite}
\fancyhead[R]{\small\color{slate}\nouppercase{\leftmark}}
\fancyfoot[L]{\scriptsize\color{slate}Independent engineering guidance; see licensed source}
\fancyfoot[R]{\scriptsize\color{slate}\thepage\ / \pageref{LastPage}}
\newtcolorbox{sourcebox}{colback=paper,colframe=navy!30,boxrule=0.45pt,arc=1.6pt,breakable}
\newtcolorbox{practicebox}{colback=warm,colframe=teal!40,boxrule=0.45pt,arc=1.6pt,breakable}
\newcommand{\SourceNote}[1]{\begin{sourcebox}\textbf{SOURCE BASIS / NORMATIVE POINTER}\quad #1\end{sourcebox}}
\newcommand{\PracticeNote}[1]{\begin{practicebox}\textbf{ENGINEERING RECOMMENDATION / NOT MISRA TEXT}\quad #1\end{practicebox}}
\newcommand{\Record}[1]{\textsf{\bfseries #1}}
\lstdefinestyle{cppsuite}{language=C++,basicstyle=\ttfamily\footnotesize,keywordstyle=\color{navy}\bfseries,commentstyle=\color{slate},stringstyle=\color{teal},breaklines=true,breakatwhitespace=true,columns=fullflexible,frame=single,framerule=0.3pt,rulecolor=\color{navy!35},backgroundcolor=\color{paper},showstringspaces=false,tabsize=2,xleftmargin=0.5em,xrightmargin=0.5em,aboveskip=4pt,belowskip=4pt}
\lstset{style=cppsuite}
\newcommand{\SuiteTitle}[3]{%
\begin{titlepage}\thispagestyle{empty}
\vspace*{1.5cm}
{\color{teal}\rule{2.4cm}{2pt}}\par\vspace{1cm}
{\Huge\bfseries\color{navy}MISRA C++:2023\\[0.4em]#1}\par\vspace{0.7cm}
{\Large\color{slate}#2}\par\vspace{1.2cm}
\begin{sourcebox}\textbf{Source baseline:} MISRA C++:2023, Guidelines for the use of C++17 in critical systems, October 2023 (user-supplied licensed copy). This is an independent companion, not an official MISRA publication, rule replacement, certified checker, or claim of project compliance.\end{sourcebox}
\vfill{\large\color{navy}Technical edition \textbullet\ 16 October 2026}\par\smallskip
{\small\color{slate}#3}\par\end{titlepage}
\tableofcontents\clearpage
}
\newcommand{\StandardsBoundary}{%
\section*{How to read normative status}\addcontentsline{toc}{section}{How to read normative status}
\textbf{Normative source:} the licensed MISRA C++:2023 guideline text, including any applicable \emph{Amplification} and \emph{Exception}. A guideline \emph{headline} alone is not a complete compliance test. The source's Section~3 defines classifications, categories, decidability, analysis scope, and applicability.\par
\textbf{Independent guidance:} all review questions, programming patterns, lifecycle processes, sample thresholds, ownership assignments, controls, and CI configurations in this suite are \emph{engineering recommendations}, not statements issued by MISRA.\par
\textbf{Critical limitation:} no individual example is asserted compliant with the whole standard. For formal compliance claims, follow the applicable edition of \emph{MISRA Compliance} as the source itself instructs, with a documented compliance plan and supporting evidence.\par
}

\begin{document}
\SuiteTitle{C++17 Example Atlas}{22 paired programming scenarios; buildable teaching fixtures}{Volume 05 / Code Examples}
\StandardsBoundary

\section{How these paired examples are scoped}
The 22 source pairs are independent, small C++17 translation units. Each noncompliant specimen is designed to illustrate an issue associated with the named guideline; the companion specimen removes that illustrated issue. They have been syntax-checked with g++ using -std=c++17, but not certified with a MISRA analyzer.

Crucially, an ``illustratively compliant snippet can still violate other guidelines. The source may have narrower amplifications, exceptions or environmental preconditions. Perform a full rule reading before using a pattern in production. Sample assets are also provided under examples/compliant and examples/noncompliant.


\section{Case 01: Control-flow reachability (Rule 0.0.1)}
\textbf{Source classification:} Required; reference printed p.~12. \textbf{Study objective:} The instruction after return cannot be reached.
\subsection*{Intentionally noncompliant pattern}
\begin{lstlisting}
// Teaching fixture: Rule 0.0.1 - Control-flow reachability
// NONCOMPLIANT for the targeted guideline only; not a whole-program compliance claim.
int compute(int n) { return n + 1; n = 3; }
\end{lstlisting}
\subsection*{Targeted correction}
\begin{lstlisting}
// Teaching fixture: Rule 0.0.1 - Control-flow reachability
// COMPLIANT for the targeted guideline only; not a whole-program compliance claim.
int compute(int n) { n += 3; return n + 1; }
\end{lstlisting}
\textbf{Verification tasks:} Identify the source amplification/exception; compile with C++17; run the MISRA rule-specific analyzer; write a one-sentence justification for the diagnosed branch; verify no new behavior regressions.

\section{Case 02: Use returned values (Rule 0.1.2)}
\textbf{Source classification:} Required; reference printed p.~20. \textbf{Study objective:} The return result is intentionally consumed; check source exceptions before dismissing a finding.
\subsection*{Intentionally noncompliant pattern}
\begin{lstlisting}
// Teaching fixture: Rule 0.1.2 - Use returned values
// NONCOMPLIANT for the targeted guideline only; not a whole-program compliance claim.
int measurement(); void sample() { measurement(); }
\end{lstlisting}
\subsection*{Targeted correction}
\begin{lstlisting}
// Teaching fixture: Rule 0.1.2 - Use returned values
// COMPLIANT for the targeted guideline only; not a whole-program compliance claim.
int measurement(); int sample() { return measurement(); }
\end{lstlisting}
\textbf{Verification tasks:} Identify the source amplification/exception; compile with C++17; run the MISRA rule-specific analyzer; write a one-sentence justification for the diagnosed branch; verify no new behavior regressions.

\section{Case 03: Parameter use (Rule 0.2.2)}
\textbf{Source classification:} Required; reference printed p.~22. \textbf{Study objective:} Removing a meaningless parameter simplifies the contract.
\subsection*{Intentionally noncompliant pattern}
\begin{lstlisting}
// Teaching fixture: Rule 0.2.2 - Parameter use
// NONCOMPLIANT for the targeted guideline only; not a whole-program compliance claim.
int square(int n, int unused) { return n * n; }
\end{lstlisting}
\subsection*{Targeted correction}
\begin{lstlisting}
// Teaching fixture: Rule 0.2.2 - Parameter use
// COMPLIANT for the targeted guideline only; not a whole-program compliance claim.
int square(int n) { return n * n; }
\end{lstlisting}
\textbf{Verification tasks:} Identify the source amplification/exception; compile with C++17; run the MISRA rule-specific analyzer; write a one-sentence justification for the diagnosed branch; verify no new behavior regressions.

\section{Case 04: Dangling pointer to automatic object (Rule 6.8.2)}
\textbf{Source classification:} Mandatory; reference printed p.~63. \textbf{Study objective:} Returning a value avoids exporting the address of a destroyed local object.
\subsection*{Intentionally noncompliant pattern}
\begin{lstlisting}
// Teaching fixture: Rule 6.8.2 - Dangling pointer to automatic object
// NONCOMPLIANT for the targeted guideline only; not a whole-program compliance claim.
int* produce() { int n = 7; return &n; }
\end{lstlisting}
\subsection*{Targeted correction}
\begin{lstlisting}
// Teaching fixture: Rule 6.8.2 - Dangling pointer to automatic object
// COMPLIANT for the targeted guideline only; not a whole-program compliance claim.
int produce() { int n = 7; return n; }
\end{lstlisting}
\textbf{Verification tasks:} Identify the source amplification/exception; compile with C++17; run the MISRA rule-specific analyzer; write a one-sentence justification for the diagnosed branch; verify no new behavior regressions.

\section{Case 05: Boolean-to-numeric conversion (Rule 7.0.1)}
\textbf{Source classification:} Required; reference printed p.~72. \textbf{Study objective:} Check the official amplification when choosing an explicit representation.
\subsection*{Intentionally noncompliant pattern}
\begin{lstlisting}
// Teaching fixture: Rule 7.0.1 - Boolean-to-numeric conversion
// NONCOMPLIANT for the targeted guideline only; not a whole-program compliance claim.
int encode(bool ready) { return ready; }
\end{lstlisting}
\subsection*{Targeted correction}
\begin{lstlisting}
// Teaching fixture: Rule 7.0.1 - Boolean-to-numeric conversion
// COMPLIANT for the targeted guideline only; not a whole-program compliance claim.
int encode(bool ready) { return ready ? 1 : 0; }
\end{lstlisting}
\textbf{Verification tasks:} Identify the source amplification/exception; compile with C++17; run the MISRA rule-specific analyzer; write a one-sentence justification for the diagnosed branch; verify no new behavior regressions.

\section{Case 06: Numeric-to-Boolean conversion (Rule 7.0.2)}
\textbf{Source classification:} Required; reference printed p.~73. \textbf{Study objective:} A meaningful predicate is clearer than an implicit numeric conversion.
\subsection*{Intentionally noncompliant pattern}
\begin{lstlisting}
// Teaching fixture: Rule 7.0.2 - Numeric-to-Boolean conversion
// NONCOMPLIANT for the targeted guideline only; not a whole-program compliance claim.
bool enabled(int value) { return value; }
\end{lstlisting}
\subsection*{Targeted correction}
\begin{lstlisting}
// Teaching fixture: Rule 7.0.2 - Numeric-to-Boolean conversion
// COMPLIANT for the targeted guideline only; not a whole-program compliance claim.
bool enabled(int value) { return value != 0; }
\end{lstlisting}
\textbf{Verification tasks:} Identify the source amplification/exception; compile with C++17; run the MISRA rule-specific analyzer; write a one-sentence justification for the diagnosed branch; verify no new behavior regressions.

\section{Case 07: Null-pointer constants (Rule 7.11.1)}
\textbf{Source classification:} Required; reference printed p.~83. \textbf{Study objective:} Use the C++17 null pointer literal.
\subsection*{Intentionally noncompliant pattern}
\begin{lstlisting}
// Teaching fixture: Rule 7.11.1 - Null-pointer constants
// NONCOMPLIANT for the targeted guideline only; not a whole-program compliance claim.
int* clear() { return 0; }
\end{lstlisting}
\subsection*{Targeted correction}
\begin{lstlisting}
// Teaching fixture: Rule 7.11.1 - Null-pointer constants
// COMPLIANT for the targeted guideline only; not a whole-program compliance claim.
int* clear() { return nullptr; }
\end{lstlisting}
\textbf{Verification tasks:} Identify the source amplification/exception; compile with C++17; run the MISRA rule-specific analyzer; write a one-sentence justification for the diagnosed branch; verify no new behavior regressions.

\section{Case 08: Avoid C-style casts (Rule 8.2.2)}
\textbf{Source classification:} Required; reference printed p.~92. \textbf{Study objective:} Named casts expose conversion intent for review.
\subsection*{Intentionally noncompliant pattern}
\begin{lstlisting}
// Teaching fixture: Rule 8.2.2 - Avoid C-style casts
// NONCOMPLIANT for the targeted guideline only; not a whole-program compliance claim.
double convert(int n) { return (double)n; }
\end{lstlisting}
\subsection*{Targeted correction}
\begin{lstlisting}
// Teaching fixture: Rule 8.2.2 - Avoid C-style casts
// COMPLIANT for the targeted guideline only; not a whole-program compliance claim.
double convert(int n) { return static_cast<double>(n); }
\end{lstlisting}
\textbf{Verification tasks:} Identify the source amplification/exception; compile with C++17; run the MISRA rule-specific analyzer; write a one-sentence justification for the diagnosed branch; verify no new behavior regressions.

\section{Case 09: Eliminate recursion (Rule 8.2.10)}
\textbf{Source classification:} Required; reference printed p.~100. \textbf{Study objective:} The loop eliminates direct recursive calls; additionally check overflow.
\subsection*{Intentionally noncompliant pattern}
\begin{lstlisting}
// Teaching fixture: Rule 8.2.10 - Eliminate recursion
// NONCOMPLIANT for the targeted guideline only; not a whole-program compliance claim.
unsigned sum(unsigned n) { return n == 0U ? 0U : n + sum(n - 1U); }
\end{lstlisting}
\subsection*{Targeted correction}
\begin{lstlisting}
// Teaching fixture: Rule 8.2.10 - Eliminate recursion
// COMPLIANT for the targeted guideline only; not a whole-program compliance claim.
unsigned sum(unsigned n) { unsigned result = 0U; for (unsigned i = 1U; i <= n; ++i) { result += i; } return result; }
\end{lstlisting}
\textbf{Verification tasks:} Identify the source amplification/exception; compile with C++17; run the MISRA rule-specific analyzer; write a one-sentence justification for the diagnosed branch; verify no new behavior regressions.

\section{Case 10: Enclose controlled statements (Rule 9.3.1)}
\textbf{Source classification:} Required; reference printed p.~114. \textbf{Study objective:} Use compound statements for condition-controlled bodies.
\subsection*{Intentionally noncompliant pattern}
\begin{lstlisting}
// Teaching fixture: Rule 9.3.1 - Enclose controlled statements
// NONCOMPLIANT for the targeted guideline only; not a whole-program compliance claim.
void set(int& n, bool on) { if (on) n = 5; }
\end{lstlisting}
\subsection*{Targeted correction}
\begin{lstlisting}
// Teaching fixture: Rule 9.3.1 - Enclose controlled statements
// COMPLIANT for the targeted guideline only; not a whole-program compliance claim.
void set(int& n, bool on) { if (on) { n = 5; } }
\end{lstlisting}
\textbf{Verification tasks:} Identify the source amplification/exception; compile with C++17; run the MISRA rule-specific analyzer; write a one-sentence justification for the diagnosed branch; verify no new behavior regressions.

\section{Case 11: Terminate conditional chains (Rule 9.4.1)}
\textbf{Source classification:} Required; reference printed p.~115. \textbf{Study objective:} An explicit final else describes the remaining case.
\subsection*{Intentionally noncompliant pattern}
\begin{lstlisting}
// Teaching fixture: Rule 9.4.1 - Terminate conditional chains
// NONCOMPLIANT for the targeted guideline only; not a whole-program compliance claim.
int rating(int n) { if (n > 0) { return 1; } else if (n < 0) { return -1; } return 0; }
\end{lstlisting}
\subsection*{Targeted correction}
\begin{lstlisting}
// Teaching fixture: Rule 9.4.1 - Terminate conditional chains
// COMPLIANT for the targeted guideline only; not a whole-program compliance claim.
int rating(int n) { if (n > 0) { return 1; } else if (n < 0) { return -1; } else { return 0; } }
\end{lstlisting}
\textbf{Verification tasks:} Identify the source amplification/exception; compile with C++17; run the MISRA rule-specific analyzer; write a one-sentence justification for the diagnosed branch; verify no new behavior regressions.

\section{Case 12: Explicit enum underlying type (Rule 10.2.1)}
\textbf{Source classification:} Required; reference printed p.~130. \textbf{Study objective:} Specify the intended underlying type and validate its target representation.
\subsection*{Intentionally noncompliant pattern}
\begin{lstlisting}
// Teaching fixture: Rule 10.2.1 - Explicit enum underlying type
// NONCOMPLIANT for the targeted guideline only; not a whole-program compliance claim.
enum class Status { pending, complete };
\end{lstlisting}
\subsection*{Targeted correction}
\begin{lstlisting}
// Teaching fixture: Rule 10.2.1 - Explicit enum underlying type
// COMPLIANT for the targeted guideline only; not a whole-program compliance claim.
enum class Status : unsigned int { pending, complete };
\end{lstlisting}
\textbf{Verification tasks:} Identify the source amplification/exception; compile with C++17; run the MISRA rule-specific analyzer; write a one-sentence justification for the diagnosed branch; verify no new behavior regressions.

\section{Case 13: Initialize before reading (Rule 11.6.2)}
\textbf{Source classification:} Mandatory; reference printed p.~138. \textbf{Study objective:} A definite initialization establishes a value on every path.
\subsection*{Intentionally noncompliant pattern}
\begin{lstlisting}
// Teaching fixture: Rule 11.6.2 - Initialize before reading
// NONCOMPLIANT for the targeted guideline only; not a whole-program compliance claim.
int dangerous() { int n; return n; }
\end{lstlisting}
\subsection*{Targeted correction}
\begin{lstlisting}
// Teaching fixture: Rule 11.6.2 - Initialize before reading
// COMPLIANT for the targeted guideline only; not a whole-program compliance claim.
int safe() { int n{0}; return n; }
\end{lstlisting}
\textbf{Verification tasks:} Identify the source amplification/exception; compile with C++17; run the MISRA rule-specific analyzer; write a one-sentence justification for the diagnosed branch; verify no new behavior regressions.

\section{Case 14: Alternatives to unions (Rule 12.3.1)}
\textbf{Source classification:} Required; reference printed p.~143. \textbf{Study objective:} std::variant carries its active alternative in C++17; assess profile and tool support.
\subsection*{Intentionally noncompliant pattern}
\begin{lstlisting}
// Teaching fixture: Rule 12.3.1 - Alternatives to unions
// NONCOMPLIANT for the targeted guideline only; not a whole-program compliance claim.
union Value { int i; float f; };
\end{lstlisting}
\subsection*{Targeted correction}
\begin{lstlisting}
// Teaching fixture: Rule 12.3.1 - Alternatives to unions
// COMPLIANT for the targeted guideline only; not a whole-program compliance claim.
#include <variant>
using Value = std::variant<int, float>;
\end{lstlisting}
\textbf{Verification tasks:} Identify the source amplification/exception; compile with C++17; run the MISRA rule-specific analyzer; write a one-sentence justification for the diagnosed branch; verify no new behavior regressions.

\section{Case 15: Prevent unintentional implicit constructors (Rule 15.1.3)}
\textbf{Source classification:} Required; reference printed p.~159. \textbf{Study objective:} An explicit constructor prevents accidental conversions at call sites.
\subsection*{Intentionally noncompliant pattern}
\begin{lstlisting}
// Teaching fixture: Rule 15.1.3 - Prevent unintentional implicit constructors
// NONCOMPLIANT for the targeted guideline only; not a whole-program compliance claim.
struct Count { Count(int n) : value(n) {} int value; };
\end{lstlisting}
\subsection*{Targeted correction}
\begin{lstlisting}
// Teaching fixture: Rule 15.1.3 - Prevent unintentional implicit constructors
// COMPLIANT for the targeted guideline only; not a whole-program compliance claim.
struct Count { explicit Count(int n) : value(n) {} int value; };
\end{lstlisting}
\textbf{Verification tasks:} Identify the source amplification/exception; compile with C++17; run the MISRA rule-specific analyzer; write a one-sentence justification for the diagnosed branch; verify no new behavior regressions.

\section{Case 16: Replace function-like macros (Rule 19.0.2)}
\textbf{Source classification:} Required; reference printed p.~182. \textbf{Study objective:} A function provides type checking and avoids multiple evaluation.
\subsection*{Intentionally noncompliant pattern}
\begin{lstlisting}
// Teaching fixture: Rule 19.0.2 - Replace function-like macros
// NONCOMPLIANT for the targeted guideline only; not a whole-program compliance claim.
#define SQUARE(x) ((x) * (x))
int square(int n) { return SQUARE(n); }
\end{lstlisting}
\subsection*{Targeted correction}
\begin{lstlisting}
// Teaching fixture: Rule 19.0.2 - Replace function-like macros
// COMPLIANT for the targeted guideline only; not a whole-program compliance claim.
constexpr int square(int n) noexcept { return n * n; }
\end{lstlisting}
\textbf{Verification tasks:} Identify the source amplification/exception; compile with C++17; run the MISRA rule-specific analyzer; write a one-sentence justification for the diagnosed branch; verify no new behavior regressions.

\section{Case 17: Automatic resource management (Rule 21.6.2)}
\textbf{Source classification:} Required; reference printed p.~200. \textbf{Study objective:} Express ownership with RAII instead of requiring callers to delete.
\subsection*{Intentionally noncompliant pattern}
\begin{lstlisting}
// Teaching fixture: Rule 21.6.2 - Automatic resource management
// NONCOMPLIANT for the targeted guideline only; not a whole-program compliance claim.
int* make_number() { return new int(4); }
\end{lstlisting}
\subsection*{Targeted correction}
\begin{lstlisting}
// Teaching fixture: Rule 21.6.2 - Automatic resource management
// COMPLIANT for the targeted guideline only; not a whole-program compliance claim.
#include <memory>
std::unique_ptr<int> make_number() { return std::make_unique<int>(4); }
\end{lstlisting}
\textbf{Verification tasks:} Identify the source amplification/exception; compile with C++17; run the MISRA rule-specific analyzer; write a one-sentence justification for the diagnosed branch; verify no new behavior regressions.

\section{Case 18: Prefer factory-created smart pointers (Rule 23.11.1)}
\textbf{Source classification:} Advisory; reference printed p.~208. \textbf{Study objective:} Factory functions encode ownership and avoid raw-pointer constructor usage.
\subsection*{Intentionally noncompliant pattern}
\begin{lstlisting}
// Teaching fixture: Rule 23.11.1 - Prefer factory-created smart pointers
// NONCOMPLIANT for the targeted guideline only; not a whole-program compliance claim.
#include <memory>
std::unique_ptr<int> make_number() { return std::unique_ptr<int>(new int(4)); }
\end{lstlisting}
\subsection*{Targeted correction}
\begin{lstlisting}
// Teaching fixture: Rule 23.11.1 - Prefer factory-created smart pointers
// COMPLIANT for the targeted guideline only; not a whole-program compliance claim.
#include <memory>
std::unique_ptr<int> make_number() { return std::make_unique<int>(4); }
\end{lstlisting}
\textbf{Verification tasks:} Identify the source amplification/exception; compile with C++17; run the MISRA rule-specific analyzer; write a one-sentence justification for the diagnosed branch; verify no new behavior regressions.

\section{Case 19: Avoid vector\textless{}bool\textgreater{} specialization (Rule 26.3.1)}
\textbf{Source classification:} Advisory; reference printed p.~215. \textbf{Study objective:} Choose storage whose element-reference semantics suit the interface.
\subsection*{Intentionally noncompliant pattern}
\begin{lstlisting}
// Teaching fixture: Rule 26.3.1 - Avoid vector<bool> specialization
// NONCOMPLIANT for the targeted guideline only; not a whole-program compliance claim.
#include <vector>
std::vector<bool> flags(8U);
\end{lstlisting}
\subsection*{Targeted correction}
\begin{lstlisting}
// Teaching fixture: Rule 26.3.1 - Avoid vector<bool> specialization
// COMPLIANT for the targeted guideline only; not a whole-program compliance claim.
#include <array>
std::array<bool, 8> flags{};
\end{lstlisting}
\textbf{Verification tasks:} Identify the source amplification/exception; compile with C++17; run the MISRA rule-specific analyzer; write a one-sentence justification for the diagnosed branch; verify no new behavior regressions.

\section{Case 20: Moved-from state (Rule 28.6.3)}
\textbf{Source classification:} Required; reference printed p.~219. \textbf{Study objective:} Re-establish a known value before accessing an object following movement.
\subsection*{Intentionally noncompliant pattern}
\begin{lstlisting}
// Teaching fixture: Rule 28.6.3 - Moved-from state
// NONCOMPLIANT for the targeted guideline only; not a whole-program compliance claim.
#include <string>
#include <utility>
std::string bad(std::string s) { auto t = std::move(s); return s + t; }
\end{lstlisting}
\subsection*{Targeted correction}
\begin{lstlisting}
// Teaching fixture: Rule 28.6.3 - Moved-from state
// COMPLIANT for the targeted guideline only; not a whole-program compliance claim.
#include <string>
#include <utility>
std::string good(std::string s) { auto t = std::move(s); s = "reset"; return s + t; }
\end{lstlisting}
\textbf{Verification tasks:} Identify the source amplification/exception; compile with C++17; run the MISRA rule-specific analyzer; write a one-sentence justification for the diagnosed branch; verify no new behavior regressions.

\section{Case 21: Use remove algorithm return value (Rule 28.6.4)}
\textbf{Source classification:} Required; reference printed p.~220. \textbf{Study objective:} The returned logical end enables erase-remove to change container size.
\subsection*{Intentionally noncompliant pattern}
\begin{lstlisting}
// Teaching fixture: Rule 28.6.4 - Use remove algorithm return value
// NONCOMPLIANT for the targeted guideline only; not a whole-program compliance claim.
#include <algorithm>
#include <vector>
void remove_zero(std::vector<int>& v) { std::remove(v.begin(), v.end(), 0); }
\end{lstlisting}
\subsection*{Targeted correction}
\begin{lstlisting}
// Teaching fixture: Rule 28.6.4 - Use remove algorithm return value
// COMPLIANT for the targeted guideline only; not a whole-program compliance claim.
#include <algorithm>
#include <vector>
void remove_zero(std::vector<int>& v) { v.erase(std::remove(v.begin(), v.end(), 0), v.end()); }
\end{lstlisting}
\textbf{Verification tasks:} Identify the source amplification/exception; compile with C++17; run the MISRA rule-specific analyzer; write a one-sentence justification for the diagnosed branch; verify no new behavior regressions.

\section{Case 22: Avoid C library I/O (Rule 30.0.1)}
\textbf{Source classification:} Required; reference printed p.~221. \textbf{Study objective:} C++ stream output can replace a C I/O call; target restrictions still apply.
\subsection*{Intentionally noncompliant pattern}
\begin{lstlisting}
// Teaching fixture: Rule 30.0.1 - Avoid C library I/O
// NONCOMPLIANT for the targeted guideline only; not a whole-program compliance claim.
#include <cstdio>
void report(int n) { std::printf("%d\n", n); }
\end{lstlisting}
\subsection*{Targeted correction}
\begin{lstlisting}
// Teaching fixture: Rule 30.0.1 - Avoid C library I/O
// COMPLIANT for the targeted guideline only; not a whole-program compliance claim.
#include <iostream>
void report(int n) { std::cout << n << "\n"; }
\end{lstlisting}
\textbf{Verification tasks:} Identify the source amplification/exception; compile with C++17; run the MISRA rule-specific analyzer; write a one-sentence justification for the diagnosed branch; verify no new behavior regressions.

\section{Running the fixture suite}
\begin{lstlisting}[language=bash,caption={Syntax-check all fixtures, not a MISRA conformance test}]
find examples -name '*.cpp' -print0 |
  xargs -0 -I{} g++ -std=c++17 -fsyntax-only '{}'
\end{lstlisting}
The fixture comparison is didactic, not an assertion that all violations are automatically detected. A classifier that reports every noncompliant fixture but also flags compliant fixtures without explanation requires deeper investigation. Evaluate each fixture against the precise official wording of its targeted guideline.


\end{document}
